\section{Introduction}
\label{sec:intro}

Autism is one diagnosis for children whose lives look very different. One child has no speech and an intellectual disability; another is in honors classes and is being treated for anxiety and depression. Much of that difference lies in the conditions that come with autism, and co-occurring conditions are a natural place to look for subtypes: groups of children who share a pattern and might share a cause, a course or a treatment.

Studies that have looked report three or four groups, usually a developmental or language group, a psychiatric or behavioral group, a small group with high rates of everything, and a large group with few conditions \citep{doshivelez2014,vargason2019,montgomery2023,litman2025}. Two problems make these findings hard to read. The first is that they are rarely checked. Of 156 autism subtyping studies reviewed by \citet{agelink2021}, 9\% repeated the analysis in an independent sample and 13\% with a second method. The second is that a clustering method finds groups whether or not groups exist. \citet{bauer2003,bauer2004} showed that mixture models recover several classes from data generated by a single continuous population, so ``the model chose $K$ classes'' is not evidence that there are $K$ kinds of children. The usual checks do not settle this either: we show below that a partition of a smooth continuum also reproduces in new samples.

This paper asks both questions of a national survey. We fit a latent class model to 18 parent-reported conditions in 8,891 autistic children aged 6 to 17 from the National Survey of Children's Health (NSCH), 2016 to 2024, and hold out the 2025 survey. The checks and their pass lines were fixed before the first model was fit. We then fit continuous latent trait models to the same data on the same likelihood scale, and we run every replication check on data simulated from each alternative, so that a passing score can be read against what a continuum with no classes would score.

Three findings follow. First, the four profiles that clinic studies report are present, as the four combinations of high and low on a developmental-language axis and a psychiatric axis, and the partition reproduces across survey years, in 2025, within age bands and under resampling. Second, a four-way cut of a smooth two-dimensional continuum reproduces about as well, while data with four real classes reproduce much better; the two-dimensional model also fits better than any class model and predicts new years as well as eight classes. Third, parent-rated severity rises along the developmental axis and does not jump at the class borders. The four profiles are reproducible regions of a map and a useful summary of it. They are not four kinds of children.

The contributions are a from-scratch implementation of the latent class model and its checks; a subtype analysis of co-occurring conditions in a nationally representative sample, with a held-out year; and a calibration of the standard replication checks against simulated data with and without classes, which applies to any subtype study that relies on them.
